%% tags: [binary search trees, time complexity]
%% source: 2023-fa-midterm_01
\begin{prob}
    Suppose a collection of unique numbers is stored in a \textbf{balanced}
    binary search tree, and that each node in the tree has been given a
    \mintinline{python}{.size} attribute which contains the number of nodes in the subtree
    rooted at that node. What is the time complexity required of an efficient
    algorithm for computing the number of elements in the collection which are
    larger than some threshold, $t$?

    \begin{soln}
        $\Theta(\log{n})$. We can consider the following algorithm: Query for $t$ in the BST, 
        and define a new variable total for tracking the result.
        For each step of the recursion, if we go to the left, add the size of the right branch 
        + 1 to the running total. We repeat the above until the query() function finishes running, 
        which will give us exactly the number of elements in the collection which are larger than 
        $t$. Since we are querying (and adding some constant time updating steps) in a balanced 
        BST, the time complexity for this question will be $\Theta(\log{n})$.
    \end{soln}

\end{prob}
